\clearpage
\section*{Problem 3}
\addcontentsline{toc}{section}{Problem 3}
\begingroup
\hypersetup{colorlinks=true,
            linkcolor=blue!50!black,
            citecolor=blue!50!black,
            urlcolor=blue!50!black}
\newcommand{\qthreeR}{\mathbb{R}}
\newcommand{\qthreeN}{\mathbb{N}}
\newcommand{\qthreeZ}{\mathbb{Z}}
\newcommand{\qthreeSnl}{S_n(\lambda)}
\newcommand{\qthreewt}{\operatorname{wt}}
\newtheorem{qthreetheorem}{Theorem}
\newtheorem{qthreelemma}{Lemma}
\newtheorem{qthreeproposition}{Proposition}
\newtheorem{qthreecorollary}{Corollary}
\theoremstyle{definition}
\newtheorem{qthreedefinition}{Definition}
\newtheorem{qthreeremark}{Remark}
\title{Solution to Problem 3: A Probabilistic Interpretation\\
for Interpolation Macdonald Polynomials}
\author{}
\date{}
\maketitle


\section*{Statement and Answer}

Let $\lambda=(\lambda_1>\cdots>\lambda_n\ge0)$ be a partition with distinct
parts that is moreover \emph{restricted}: it has a unique part of size $0$ and no
part of size $1$. Let $\qthreeSnl$ be the orbit of $\lambda$ under permutation of
parts (equivalently, the compositions that rearrange $\lambda$). Let
$F^*_\mu(x_1,\dots,x_n;q,t)$ and $P^*_\lambda(x_1,\dots,x_n;q,t)$ be the
interpolation ASEP polynomial and interpolation Macdonald polynomial.

\medskip
\noindent\textbf{Theorem.} \emph{There is a nontrivial Markov chain on $\qthreeSnl$ ---
the \emph{interpolation $t$-Push TASEP} --- whose stationary distribution is}
\[
\pi^*_\lambda(\mu)\;=\;\frac{F^*_\mu(x_1,\dots,x_n;\,q=1,\,t)}{P^*_\lambda(x_1,\dots,x_n;\,q=1,\,t)},
\qquad \mu\in\qthreeSnl,
\]
\emph{and whose transition probabilities are not described using the
polynomials $F^*_\mu$.}

\medskip
\noindent The answer is \textbf{YES}. We caution at the outset that the chain is
\emph{not} the ordinary multispecies ASEP: that chain has the ordinary (not
interpolation) Macdonald polynomials as its stationary weights (Ayyer--Martin--Williams),
and substituting it here would solve a different, already-known problem. The
correct object is a $q=1$ \emph{push}-TASEP with an interpolation correction
encoded by a ringing ``bell''.

\section*{Probabilities and notation}

Throughout we work at $q=1$ and write $P^*_\lambda=P^*_\lambda(x;1,t)$, etc.
For $1\le k\le n$ put
\[
p_k:=\frac{t^{-n+1}(1-t)}{x_k-t^{-n+2}},\qquad
q_k:=\frac{(1-t)\,x_k}{x_k-t^{-n+2}} .
\]
For $0<t<1$ and $x_i>t^{-n+1}$ these lie in $[0,1]$ and serve as
skipping/displacement probabilities. We use $e^*_k(x;t)$ for the interpolation
elementary symmetric polynomials and the factorization (valid at $q=1$
\cite{Dolega,BDW25a})
\begin{equation}\label{eq:factor}
P^*_\lambda(x;1,t)=\prod_{1\le i\le\lambda_1}e^*_{\lambda'_i}(x;t),
\end{equation}
$\lambda'$ the conjugate partition.

We record two elementary identities used repeatedly below; they are the
algebraic content of the correspondence in Proposition~\ref{prop:step2}. Since
$t^{-n+2}-t^{-n+1}(1-t)=t^{-n+1}\bigl(t-(1-t)\bigr)$ is \emph{not} what we want,
we compute directly:
\begin{align}
\bigl(x_k-t^{-n+2}\bigr)-t^{-n+1}(1-t)
   &=x_k-t^{-n+2}-t^{-n+1}+t^{-n+2}=x_k-t^{-n+1},
   \label{eq:idp}\\
\bigl(x_k-t^{-n+2}\bigr)-(1-t)\,x_k
   &=t\,x_k-t^{-n+2}=t\bigl(x_k-t^{-n+1}\bigr).
   \label{eq:idq}
\end{align}
Consequently, dividing \eqref{eq:idp} and \eqref{eq:idq} by $x_k-t^{-n+2}$,
\begin{equation}\label{eq:skipprobs}
1-p_k=\frac{x_k-t^{-n+1}}{x_k-t^{-n+2}},\qquad
1-q_k=\frac{t\,\bigl(x_k-t^{-n+1}\bigr)}{x_k-t^{-n+2}} .
\end{equation}

\section*{The interpolation $t$-Push TASEP}

\begin{qthreedefinition}\label{def:chain}
Fix $\lambda$ with $\lambda_n=0$. States of the chain are configurations of
particles on a ring of $n$ sites labelled by the parts of $\lambda$: state $\eta$
places the particle labelled $\eta_j$ at site $j$ (label $0$ is a vacancy). One
step from $\eta$ proceeds in three stages.

\smallskip
\noindent\textbf{(Step 0) The bell.} The bell at site $j$ rings with probability
\[
P_j=\frac{\prod_{k<j}\bigl(x_k-t^{-n+2}\bigr)\,\prod_{k>j}\bigl(x_k-t^{-n+1}\bigr)}
{e^*_{n-1}(x;t)},
\quad
e^*_{n-1}(x;t)=\sum_{j=1}^n\prod_{k<j}\bigl(x_k-t^{-n+2}\bigr)\prod_{k>j}\bigl(x_k-t^{-n+1}\bigr).
\]

\smallskip
\noindent\textbf{(Step 1) Push-TASEP move.} The particle at site $j$, say with
label $a$, is activated and travels clockwise by the $t$-Push TASEP rule: if
there are $m$ weaker particles in the system (labels $<a$, vacancies counting as
label $0$), the active particle moves to the $k$-th such weaker particle with
probability $t^{k-1}/[m]_t$. If that site held a positive label, that particle
becomes active and repeats; the cascade ends when the active particle reaches a
vacancy. Afterwards site $j$ is vacant; we treat this vacancy as a new particle
labelled $a:=0$.

\smallskip
\noindent\textbf{(Step 2) Vacancy sweep.} The label-$0$ particle goes to site $1$
and travels clockwise. Reaching a site $k$ ($1\le k\le j-1$) holding a label
$b\ge0$, it skips that site with probability $1-p_k$ if $b\ge a$ and $1-q_k$ if
$b<a$; otherwise it settles there, activating/displacing the resident, which then
continues clockwise toward site $j$ by the same rule. The active particle stops
once it displaces another or arrives at site $j$, where it settles.
\end{qthreedefinition}

We denote the transition probabilities of Steps 1 and 2 by $P^{(1)}_j$ and
$P^{(2)}_j$, so that for $\mu,\nu\in\qthreeSnl$,
\[
P(\mu,\nu)=\sum_{1\le j\le n}P_j\sum_{\rho\in\qthreeSnl:\rho_j=0}
P^{(1)}_j(\mu,\rho)\,P^{(2)}_j(\rho,\nu).
\]

\begin{qthreeremark}[nontriviality]
The rates $P_j,p_k,q_k$ are explicit rational functions of $x_1,\dots,x_n$ and
$t$ only. They do \emph{not} reference $F^*_\mu$, so the chain is nontrivial in
the required sense (unlike a Metropolis--Hastings chain built from the target
distribution).
\end{qthreeremark}

\begin{qthreetheorem}\label{thm:main}
In the interpolation $t$-Push TASEP with content $\lambda$, the stationary
probability of $\mu\in\qthreeSnl$ is $\pi^*_\lambda(\mu)=F^*_\mu(x;1,t)/P^*_\lambda(x;1,t)$.
\end{qthreetheorem}

\section*{Two-line queues and the proof}

We use classical two-line queues \cite{CMW22} and signed two-line queues
\cite{BDW25a}, specialized to $q=1$. Let $\mathcal Q^\eta_\kappa$ be the classical
two-line queues with top row $\eta$ and bottom row $\kappa$, with pair-weight
generating function $a^\eta_\kappa=\sum_{Q\in\mathcal Q^\eta_\kappa}\qthreewt_{\mathrm{pair}}(Q)$.
Let $\mathcal G^\alpha_\mu$ be the signed two-line queues with top row $\alpha\in\qthreeZ^n$
and bottom row $\mu$, with $b^\alpha_\mu=\sum_{Q\in\mathcal G^\alpha_\mu}\qthreewt_{\mathrm{pair}}(Q)$,
and full weight $\qthreewt(Q)=\qthreewt_{\mathrm{pair}}(Q)\qthreewt_{\mathrm{ball}}(Q)$ where
\begin{equation}\label{eq:wta}
\qthreewt_\alpha\,b^\alpha_\mu=\sum_{Q\in\mathcal G^\alpha_\mu}\qthreewt(Q),
\quad \qthreewt_\alpha:=\prod_{k:\alpha_k>0}x_k\prod_{k:\alpha_k<0}\bigl(-t^{-(n-1)}\bigr).
\end{equation}

\begin{qthreedefinition}\label{def:unsigned}
For a signed two-line queue $Q\in\mathcal G^\alpha_\mu$ let $\overline Q$ be its
unsigned version (forget signs on the top row); its top composition is
$\|\alpha\|=(|\alpha_1|,\dots,|\alpha_n|)$. Let $\mathcal G^\kappa_\mu$ be the set of
unsigned objects so obtained, with $\|\alpha\|=\kappa$. For $\overline Q\in\mathcal G^\kappa_\mu$
assign a nontrivial pairing $p$ the weight $(1-t)t^{\mathrm{skip}(p)}$, where
$\mathrm{skip}(p)$ is the number of weaker balls (labels strictly smaller than the
moving label, vacancies included) over which $p$ jumps, and a top-row
ball $B$ in column $k$ labelled $a>0$ with the ball below labelled $b$ the weight
\[
\qthreewt(B)=\begin{cases}x_k-t^{-(n-1)} & b=a,\\ x_k & b>a,\\ t^{-(n-1)} & b<a.\end{cases}
\]
\end{qthreedefinition}

\begin{qthreelemma}\label{lem:signforget}
For $\lambda$ with distinct parts and $\kappa,\mu\in\qthreeSnl$, and $\overline Q\in\mathcal G^\kappa_\mu$,
$\qthreewt(\overline Q)=\sum_Q\qthreewt(Q)$, the sum over signed $Q$ forgetting to $\overline Q$.
\end{qthreelemma}

\begin{proof}
Enumerate the admissible sign choices for each top-row ball $B$ labelled $a>0$:
a $-$ sign is forced when the ball below is a vacancy or labelled $b<a$; a $+$
sign is forced when $b>a$; both are allowed when $b=a$. One checks this matches
\eqref{eq:wta}: a $-$ sign multiplies the ball weight by $-1$ in passing from
$Q$ to $\overline Q$, but the weight of the pairing attached to $B$ is also
multiplied by $-1$, so signs cancel consistently.
\end{proof}

For $\kappa\in\qthreeSnl$ set $c^\kappa_\mu:=\sum_{\alpha:\|\alpha\|=\kappa}\qthreewt_\alpha\,b^\alpha_\mu$.
Combining \eqref{eq:wta} and Lemma~\ref{lem:signforget}:

\begin{qthreelemma}\label{lem:ckm}
$c^\kappa_\mu=\sum_{\overline Q\in\mathcal G^\kappa_\mu}\qthreewt(\overline Q)$. Since $\lambda$ has distinct
parts, $\mathcal G^\kappa_\mu$ is empty or a single element.
\end{qthreelemma}

\subsection*{The closing recursion (where the hypotheses on $\lambda$ enter)}

Fix a weakly order-preserving $\phi:\qthreeN\to\qthreeN$ and partitions $\lambda,\kappa$ with
$\phi(\lambda)=\kappa$. For $\eta\in S_n(\kappa)$ set
$G^*_\eta(x;t):=\sum_{\rho\in\qthreeSnl:\,\phi(\rho)=\eta}F^*_\rho(x;1,t)$.

\begin{qthreetheorem}[decrement projection]\label{thm:G}
Let $\lambda$ be restricted, $\phi(i)=\max(i-1,0)$, and $\kappa=\phi(\lambda)=\lambda^-$.
Then for all $\eta\in S_n(\kappa)$, at $q=1$,
\[
\frac{G^*_\eta(x;t)}{P^*_\lambda(x;1,t)}=\frac{F^*_\eta(x;1,t)}{P^*_\kappa(x;1,t)} .
\]
\end{qthreetheorem}

This is the only case of the projection identity used in the proof of
Theorem~\ref{thm:main}: it is exactly the column-removal recursion of
Corollary~\ref{cor:rec}. We prove it below by a \emph{complete}
weight-preserving bijection (Lemma~\ref{lem:decrement}); the restriction on
$\lambda$ (no part of size $1$) is what makes the decrement \emph{strictly}
order-preserving on the occurring labels (Lemma~\ref{lem:strict}), so that no
column case is ever altered and the bijection is unconditional.

\begin{qthreeremark}[scope: the general projection identity]\label{rmk:general}
For a general weakly order-preserving deflating $\phi$ with $\phi(\lambda)=\kappa$
one expects the full analogue of the Ayyer--Martin--Williams projection identity
\eqref{eq:AMWproj},
\[
\frac{G^*_\eta(x;t)}{P^*_\lambda(x;1,t)}=\frac{F^*_\eta(x;1,t)}{P^*_\kappa(x;1,t)},
\qquad \eta\in S_n(\kappa),
\]
to persist for the interpolation polynomials. We state this as
Proposition~\ref{prop:general} below and give the part of its proof that is
genuinely established (the reduction to elementary moves and the complete proof of
the \emph{gap-decrement} move), together with a precise account of the one
remaining ingredient --- the \emph{merge} move --- which we do \emph{not} prove
here and which is \emph{not} needed for any other result in this paper. We are
careful to label exactly what is proven and what is not, rather than asserting the
merge identity. We first record the combinatorial model, the precise statement of
the result being lifted, and the new ingredient.
\end{qthreeremark}

\paragraph{The combinatorial model.}
We use the combinatorial formula for the interpolation ASEP polynomials at $q=1$
established in \cite[Thm.~1.15]{BDW25a}; it is the same formula underlying the
expansion used in the proof of Theorem~\ref{thm:main}. For a partition
$\lambda=(\lambda_1>\cdots>\lambda_n\ge0)$ with $L:=\lambda_1$, a
\emph{signed multiline queue} (sMLQ) of content $\lambda$ is a stack of $L$ rows
indexed $r=1,\dots,L$ on the ring of $n$ sites, where row $r$ carries the balls
of every label $a$ with $a\ge r$; equivalently the \emph{column} of a label-value
$v$ occupies rows $1,\dots,v$. The bottom row (row $1$) reads off a composition
$\nu\in S_n(\lambda)$, and \cite[Thm.~1.15]{BDW25a} gives
\begin{equation}\label{eq:Fformula}
F^*_\nu(x;1,t)=\sum_{Q\in\mathrm{sMLQ}(\lambda):\ \mathrm{bot}(Q)=\nu}\qthreewt(Q),
\end{equation}
where $\qthreewt(Q)$ is the product, over all columns of all rows, of the ball
weights of Definition~\ref{def:unsigned}
\[
\qthreewt(B)=\begin{cases}x_k-t^{-(n-1)} & b=a,\\ x_k & b>a,\\ t^{-(n-1)} & b<a,\end{cases}
\]
($k$ the site, $a$ the label of $B$, $b$ the label directly below $B$ in the next
lower row, with the convention $b=a$ in the bottom row), times the pairing weights
$(1-t)t^{\mathrm{skip}(p)}$ over the nontrivial pairings $p$ linking each row to
the one below. Two facts about \eqref{eq:Fformula}, immediate from
Definition~\ref{def:unsigned}, are used repeatedly:
\begin{itemize}
\item[(W1)] \emph{(case-locality)} the ball weight of $B$ depends only on its site
$k$ and on the three-way comparison of its label $a$ with the label $b$ directly
below it ($a=b$, $a<b$, $a>b$);
\item[(W2)] \emph{(skip-locality)} $\mathrm{skip}(p)$ counts the balls in the
lower row that are \emph{strictly weaker} (smaller label, vacancies counting as
label $0$) than the moving label of $p$, so it too depends only on the relative
order of labels.
\end{itemize}
The row factorization of $P^*$ at $q=1$, \eqref{eq:factor}, reads here
\begin{equation}\label{eq:Prow}
P^*_\lambda(x;1,t)=\prod_{r=1}^{L}e^*_{\lambda'_r}(x;t),\qquad
\lambda'_r=\#\{i:\lambda_i\ge r\}=\#\{\text{balls in row }r\},
\end{equation}
and $e^*_m(x;t)$ is the total weight of a single $m$-ball row over the prescribed
level, the case $m=n-1$ being the formula displayed in Definition~\ref{def:chain};
the single-row generating identity (the $q=1$ Pieri/row rule) is
\cite[Lem.~5.4]{AMW25} in the ordinary case and \cite[Lem.~5.6]{BDW25a} in the
interpolation case.

\paragraph{The result being lifted, and the new ingredient.}
For \emph{ordinary} ($q=1$) ASEP polynomials $F_\mu$ and Macdonald polynomials
$P_\lambda$, the projection identity is the color-merging/projection of multiline
queues of Ayyer--Martin--Williams \cite[Prop.~6.1]{AMW25} (in the lineage of
\cite[\S5]{CMW22}): for weakly order-preserving $\phi$ with $\phi(\lambda)=\kappa$,
\begin{equation}\label{eq:AMWproj}
\sum_{\rho:\ \phi(\rho)=\eta}F_\rho(x;1,t)=\frac{P_\lambda(x;1,t)}{P_\kappa(x;1,t)}\,F_\eta(x;1,t),
\qquad \eta\in S_n(\kappa).
\end{equation}
Its mechanism: applying $\phi$ to all labels of a multiline queue of content
$\lambda$ produces one of content $\kappa$; for the ordinary (trivial/$x_k$) ball
weights this is weight-preserving, and the fiber over a fixed $\kappa$-queue with
bottom $\eta$ has total weight $P_\lambda/P_\kappa$. We use \cite[Thm.~1.1]{AS19}
only to guarantee that the $q=1$ specialization of the multiline-queue weights is
well defined (rational in $x,t$, no residual $q$-dependence), which is what makes
the projection map and its fibers meaningful at $q=1$.

The genuinely new feature of the \emph{interpolation} weights is that the ball
weights $x_k-t^{-(n-1)},\,x_k,\,t^{-(n-1)}$ are \emph{case-sensitive} by (W1):
color merging can change the comparison $a$ vs.\ $b$ in a column, hence its weight,
so \eqref{eq:AMWproj} is not formally inherited. Whether the interpolation
projection nonetheless holds after summing the fiber is settled below \emph{only}
for moves that do not merge two present values
(Proposition~\ref{prop:general}\,\textup{(B)}); the merging case is isolated as the
unresolved ingredient \eqref{eq:mergewt}. The decrement case used by
Theorem~\ref{thm:main} avoids the issue entirely, never merging:

\begin{qthreelemma}\label{lem:strict}
Let $\phi(i)=\max(i-1,0)$ and suppose $\lambda$ has no part equal to $1$, so all
labels lie in $V:=\{0,2,3,4,\dots\}$. Then $\phi$ is \emph{strictly} increasing on
$V$: for $a,b\in V$, $\,a<b\iff\phi(a)<\phi(b)\,$ and $\,a=b\iff\phi(a)=\phi(b)$.
\end{qthreelemma}

\begin{proof}
If $a=b$ then $\phi(a)=\phi(b)$. If $0=a<b$ then $b\ge2$, so
$\phi(a)=0<1\le b-1=\phi(b)$. If $2\le a<b$ then $\phi(a)=a-1<b-1=\phi(b)$. Thus
$a<b\Rightarrow\phi(a)<\phi(b)$, and the equivalences follow since $\phi$ is
nondecreasing. The clamping at $0$ never merges two distinct values of $V$
precisely because $1\notin V$: the only value of $V$ sent to $0$ is $0$ itself.
\end{proof}

\begin{qthreelemma}[decrement projection]\label{lem:decrement}
Let $\lambda$ be restricted, $\phi(i)=\max(i-1,0)$, $\kappa=\phi(\lambda)=\lambda^-$,
and $k=\lambda'_1$ the number of nonzero parts of $\lambda$. Then for every
$\eta\in S_n(\kappa)$ there is exactly one $\rho\in S_n(\lambda)$ with
$\phi(\rho)=\eta$, namely $\rho$ obtained from $\eta$ by replacing each part $\ge1$
by that part $+1$, and
\[
F^*_\rho(x;1,t)=e^*_k(x;t)\,F^*_\eta(x;1,t).
\]
\end{qthreelemma}

\begin{proof}
\emph{Uniqueness of the preimage.} Since $1\notin\{\lambda_i\}$, Lemma~\ref{lem:strict}
makes $\phi$ a bijection $V\to\qthreeN$, with inverse $0\mapsto0$ and $m\mapsto m+1$
for $m\ge1$. Applying $\phi^{-1}$ entrywise to $\eta$ yields the unique
$\rho\in S_n(\lambda)$ with $\phi(\rho)=\eta$; hence $G^*_\eta=F^*_\rho$.

\emph{Bottom-row deletion realises the decrement.} Recall the model: row $r$
($1\le r\le L$) carries exactly the balls of label $\ge r$, so a label-value $v$
occupies rows $1,\dots,v$ and row $r$ has $\lambda'_r=\#\{i:\lambda_i\ge r\}$ balls.
Relabelling every ball of label $a$ by $\phi(a)=\max(a-1,0)$ defines
\[
\Psi:\mathrm{sMLQ}(\lambda)\longrightarrow\mathrm{sMLQ}(\kappa).
\]
Because $\lambda$ is restricted, all nonzero labels are $\ge2$, so by
Lemma~\ref{lem:strict} a ball lies in \emph{new} row $r$ (i.e.\ $\phi(a)\ge r$) iff
it lay in \emph{old} row $r+1$ (i.e.\ $a\ge r+1$):
\begin{equation}\label{eq:rowshift}
(\text{new row }r)=(\text{old row }r+1),\qquad 1\le r\le L-1,
\end{equation}
the labels relabelled by the order-isomorphism $\phi$. Thus \emph{$\Psi$ deletes the
old bottom row $($row $1$, of $\lambda'_1=k=n-1$ balls$)$ and shifts every higher row
down by one}: the bottom row $\nu$ becomes $\phi(\nu)=\eta$, content $\lambda$ becomes
$\kappa=\lambda^-$, and $(\lambda^-)'_r=\lambda'_{r+1}$ matches \eqref{eq:Prow}.
(The originally tempting description --- delete ``the top ball of each column'' ---
is \emph{not} a single level: those balls sit at the distinct rows
$\lambda_1>\dots>\lambda_k$, whereas the decrement is genuinely the deletion of the
single bottom row, which is what \eqref{eq:rowshift} establishes and what makes the
single-row rule below applicable.)

\emph{(a) The surviving rows keep their weight.} Since $\phi$ is an order-isomorphism
on the occurring labels (Lemma~\ref{lem:strict}), every three-way comparison (W1) and
every skip-count (W2) \emph{among the surviving rows} $r\ge2$ of $Q$ --- which are
carried verbatim, with relabelled labels, to rows $r\ge1$ of $\Psi(Q)$ by
\eqref{eq:rowshift} --- is preserved. Hence the contribution of rows $\ge3$ of $Q$
(together with all pairings among rows $\ge2$) equals the contribution of rows
$\ge2$ of $\Psi(Q)$; the only data of $Q$ not already present in $\Psi(Q)$ are the
deleted bottom row and its interface (ball weights and pairings) to old row $2$.
Write $\qthreewt_{\mathrm{ins}}(Q):=\qthreewt(Q)/\qthreewt(\Psi(Q))$ for this
leftover factor --- the weight that the deleted bottom row and its interface
contribute on top of $\Psi(Q)$.

\emph{(b) The fiber is a single bottom-row insertion of weight $e^*_k$, independent of
the image.} Fix $\overline Q\in\mathrm{sMLQ}(\kappa)$ with $\mathrm{bot}(\overline Q)=\eta$.
By \eqref{eq:rowshift} the preimage $\Psi^{-1}(\overline Q)$ is exactly the set of
sMLQs obtained by inserting \emph{one} new bottom row of $k=n-1$ balls --- carrying the
$n-1$ distinct nonzero labels of $\lambda$, on the sites prescribed by the unique
preimage $\rho$ --- beneath the fixed configuration $\overline Q$, with all admissible
pairings from $\overline Q$'s bottom row down to the inserted row. This is precisely the
input of the $q=1$ interpolation single-row $($Pieri$)$ rule
\cite[Lem.~5.6]{BDW25a} $(\cite[Lem.~5.4]{AMW25}$ in the ordinary case$)$: the total
weight of one fresh $m$-ball row inserted adjacent to a fixed configuration, summed
over the inserted row's ball cases and pairings, equals the interpolation elementary
symmetric polynomial $e^*_m(x;t)$, and --- this is the substantive content of that
lemma --- the value is \emph{independent of the adjacent configuration}. With $m=k$,
\begin{equation}\label{eq:fibertop}
\sum_{Q\in\Psi^{-1}(\overline Q)}\qthreewt_{\mathrm{ins}}(Q)=e^*_k(x;t)
\qquad\text{for every }\overline Q\in\mathrm{sMLQ}(\kappa).
\end{equation}
This is the only point at which the row rule \eqref{eq:Prow} is used, and it is used
for a genuine single row (the deleted bottom row), not for a staircase. The two
hypotheses on $\lambda$ place us exactly in its scope: \emph{no part of size $1$}
makes $\phi$ an order-isomorphism, so that $\Psi$ is the clean bottom-row deletion
\eqref{eq:rowshift} and no surviving column case is disturbed (Lemma~\ref{lem:strict});
without it two labels would merge, a comparison $a>b$ could collapse to $a=b$, and one
would be in the unproven merge regime of Proposition~\ref{prop:general}\,\textup{(C)};
and \emph{distinct parts} make the $k$ inserted labels pairwise distinct, so that the
inserted row is a genuine $(n-1)$-ball row of the model.

\emph{The value $e^*_k$, pinned non-circularly.} Equation \eqref{eq:fibertop} already
follows from the cited single-row rule. Independently --- granting only the
\emph{qualitative} half of that rule, that the fiber weight is \emph{some} constant
$E_k$ independent of $\overline Q$ --- the value is forced by the symmetrization
identity $\sum_{\mu\in S_n(\lambda)}F^*_\mu(x;1,t)=P^*_\lambda(x;1,t)$ (the $q=1$ case
of the standard fact that the interpolation ASEP polynomials of content $\lambda$ sum
to the interpolation Macdonald polynomial; it does \emph{not} use
Corollary~\ref{cor:rec}, so the check is non-circular). Indeed parts (a)--(b) then give
$F^*_\rho=E_k\,F^*_\eta$ for the unique preimage $\rho$ of each $\eta\in S_n(\kappa)$;
summing over $\eta\in S_n(\kappa)$ and using that $\Psi$ matches the orbits
$S_n(\lambda)\leftrightarrow S_n(\kappa)$ bijectively (Lemma~\ref{lem:strict}),
\[
P^*_\lambda=\sum_{\rho\in S_n(\lambda)}F^*_\rho
=E_k\sum_{\eta\in S_n(\kappa)}F^*_\eta=E_k\,P^*_\kappa,
\]
so $E_k=P^*_\lambda/P^*_\kappa=e^*_{\lambda'_1}(x;t)=e^*_k(x;t)$, the last step by the
telescoping of \eqref{eq:Prow} (using $\lambda'_1=k$). Thus the constant fiber value is
$e^*_k$, agreeing with \eqref{eq:fibertop}.

Combining (a) and (b), every $Q$ with $\mathrm{bot}(Q)=\rho$ maps to $\Psi(Q)$ with
$\mathrm{bot}(\Psi(Q))=\phi(\rho)=\eta$, and using
$\qthreewt(Q)=\qthreewt(\Psi(Q))\,\qthreewt_{\mathrm{ins}}(Q)$,
\[
\sum_{Q:\ \mathrm{bot}=\rho}\qthreewt(Q)
=\sum_{\overline Q:\ \mathrm{bot}=\eta}\ \qthreewt(\overline Q)\!\!\sum_{Q\in\Psi^{-1}(\overline Q)}\!\!\qthreewt_{\mathrm{ins}}(Q)
=e^*_k(x;t)\!\!\sum_{\overline Q:\ \mathrm{bot}=\eta}\!\!\qthreewt(\overline Q).
\]
By \eqref{eq:Fformula} the left side is $F^*_\rho$ (as $\rho$ is the unique preimage
of $\eta$) and the right side is $e^*_k\,F^*_\eta$. This proves the lemma.
\end{proof}

\begin{proof}[Proof of Theorem~\ref{thm:G}]
Here $\phi(i)=\max(i-1,0)$ and $\lambda$ is restricted, so $\kappa=\lambda^-$. By
Lemma~\ref{lem:decrement} the preimage of $\eta$ is the unique $\rho\in S_n(\lambda)$
with $\phi(\rho)=\eta$, whence $G^*_\eta=F^*_\rho=e^*_k\,F^*_\eta$. The factorization
\eqref{eq:Prow} telescopes under decrement: $(\lambda^-)'_r=\lambda'_{r+1}$, so
\[
\frac{P^*_\lambda}{P^*_\kappa}
=\frac{\prod_{r\ge1}e^*_{\lambda'_r}}{\prod_{r\ge1}e^*_{\lambda'_{r+1}}}
=e^*_{\lambda'_1}=e^*_k ,
\]
the last equality because $\lambda'_1=k$ is the number of nonzero parts. Dividing,
$G^*_\eta/P^*_\lambda=F^*_\eta/P^*_\kappa$.
\end{proof}

\bigskip

We now turn to the general projection identity announced in
Remark~\ref{rmk:general}. We state it as a proposition and prove exactly as much as
we can establish rigorously, marking precisely the one step that we do not prove.

\begin{qthreeproposition}[general projection identity; merge move unproven]\label{prop:general}
Let $\phi$ be weakly order-preserving and deflating with $\phi(\lambda)=\kappa$,
$\lambda$ with distinct parts. Then \eqref{eq:fiberwt} below --- equivalently the
identity $G^*_\eta/P^*_\lambda=F^*_\eta/P^*_\kappa$ for all $\eta\in S_n(\kappa)$ ---
reduces to a single-move statement for each elementary move \textup{(i)},
\textup{(ii)} of the factorization established in part \textup{(A)} of the proof. We
prove the reduction \textup{(A)} and the \emph{gap-decrement} move \textup{(i)}
\textup{(B)} in full. The \emph{merge} move \textup{(ii)} \textup{(C)} is the
genuinely new interpolation ingredient; we give the precise statement that must hold
and the structural reason it is expected, but we do \emph{not} prove it here. The
identity therefore holds for every $\phi$ whose factorization uses only move
\textup{(i)} (in particular for the decrement of Theorem~\ref{thm:G}, which uses no
move \textup{(ii)}); for general $\phi$ it is conditional on \textup{(C)}.
\end{qthreeproposition}

\begin{proof}
\emph{Setup.} Applying $\phi$ entrywise to all labels of an sMLQ defines
\[
\Phi:\mathrm{sMLQ}(\lambda)\to\mathrm{sMLQ}(\kappa),\qquad
\mathrm{bot}(\Phi(Q))=\phi(\mathrm{bot}(Q)),
\]
which carries $\{Q:\mathrm{bot}(Q)=\rho,\ \phi(\rho)=\eta\}$ into
$\{\overline Q:\mathrm{bot}(\overline Q)=\eta\}$. Summing the per-fiber identity
\begin{equation}\label{eq:fiberwt}
\sum_{Q\in\Phi^{-1}(\overline Q)}\qthreewt(Q)
=\frac{P^*_\lambda(x;1,t)}{P^*_\kappa(x;1,t)}\,\qthreewt(\overline Q),
\qquad \overline Q\in\mathrm{sMLQ}(\kappa),
\end{equation}
over $\overline Q$ with $\mathrm{bot}=\eta$ and applying \eqref{eq:Fformula} on both
sides yields $G^*_\eta=(P^*_\lambda/P^*_\kappa)F^*_\eta$. So it suffices to prove
\eqref{eq:fiberwt}.

\emph{(A) Factorization into elementary moves.} Since $\kappa=\phi(\lambda)$ has each
part $\le$ the corresponding part of $\lambda$, $\phi$ is deflating: $\phi(v)\le v$
for every present value $v$. On the set $U$ of present label values consider the two
elementary order-preserving moves:
\begin{itemize}
\item[(i)] \emph{(gap decrement)} for a present value $v$ with $v\ge1$ and
$v-1\notin U$, replace every label $v$ by $v-1$; this is injective and strictly
order-preserving on $U$ (the slot $v-1$ is empty, so no two distinct present values
collide, by the case analysis of Lemma~\ref{lem:strict});
\item[(ii)] \emph{(merge)} for present $v-1<v$ both in $U$, identify them, replacing
every label $v$ by $v-1$.
\end{itemize}
\emph{Every deflating order-preserving $\phi$ is a finite composition of moves
\textup{(i)},\textup{(ii)}.} Run the following procedure on a current labelling $g$
(initially $g=\mathrm{id}_U$, target $\phi$). While some present value $v$ has
$g(v)>\phi(v)$, pick the smallest such $v$ and lower its block by one: apply move
\textup{(i)} if $g(v)-1\notin g(U)$, and move \textup{(ii)} if $g(v)-1\in g(U)$. The
latter is legitimate: $v$ being the smallest block currently above its target and
$\phi$ being order-preserving, the value currently occupying slot $g(v)-1$ has the
same final target as $v$ (both equal $\phi(v)=\phi(v-1)$ once they meet), so the two
are slated to merge; hence move \textup{(ii)} applied now is a sub-step of $\phi$.
Each application strictly decreases the nonnegative integer
$\sum_{w\in U}\bigl(g(w)-\phi(w)\bigr)$ by exactly the size of the block moved, so the
procedure terminates, necessarily at $g=\phi$ (it halts only when no block exceeds its
target, i.e.\ $g=\phi$). This proves the factorization for \emph{every} deflating
order-preserving $\phi$; the explicit machine check on the value sets
$\{0,2,3\},\{0,2,4\},\{0,2,3,4\},\{0,2,3,5\}$ is merely a sanity confirmation, not the
proof. Because $\Phi$ is the entrywise composite of the corresponding single-move
relabellings, fiber weights multiply along the composition, and \eqref{eq:Prow} shows
the product of the per-move row-factor ratios telescopes to $P^*_\lambda/P^*_\kappa$.
Thus \eqref{eq:fiberwt} for $\phi$ follows once it is proved for a single move, with
$P^*_\lambda/P^*_\kappa$ replaced by the relevant ratio of $P^*$ row factors.

\emph{(B) The gap-decrement move \textup{(i)} (proved).} Here $v$ is not merged with
any value (the slot $v-1$ is empty), so by Lemma~\ref{lem:strict} the three-way
comparison in every column is preserved across all levels, and no ball case changes.
The single relabelling $v\mapsto v-1$ then frees level $v$ to be deleted exactly as in
Lemma~\ref{lem:decrement} applied to the single row $r=v$: the fiber over the image
queue is the set of ways to install one fresh level of
$\lambda'_v=\#\{\text{balls of value }\ge v\}$ balls over the (unchanged) row below,
and the verbatim argument of Lemma~\ref{lem:decrement}(a)--(b) --- which uses only
case-locality (W1), skip-locality (W2), and the single-row generating identity
\eqref{eq:Prow} for the base instance --- gives fiber weight $e^*_{\lambda'_v}$. Since
the image retains a single row of the same size $\lambda'_v$ at level $v-1$, this is
exactly the row-factor ratio for move \textup{(i)}. Hence \eqref{eq:fiberwt} holds for
every move \textup{(i)}.

\emph{(C) The merge move \textup{(ii)} (statement and obstruction; not proved here).}
Here $v-1,v\in U$ are identified, level $v$ is deleted, and levels $v-1,v$ collapse to
a single image level at $\phi(v)=v-1$ carrying $\kappa'_{\phi(v)}=\lambda'_{v-1}$ balls.
The required single-move instance of \eqref{eq:fiberwt} reads
\begin{equation}\label{eq:mergewt}
\sum_{Q\in\Phi^{-1}(\overline Q)}\qthreewt_{\{v-1,v\}}(Q)
=\frac{e^*_{\lambda'_{v-1}}\,e^*_{\lambda'_v}}{e^*_{\kappa'_{\phi(v)}}}\,
\qthreewt_{\phi(v)}(\overline Q)
=e^*_{\lambda'_v}\,\qthreewt_{\phi(v)}(\overline Q),
\end{equation}
where $\qthreewt_{\{v-1,v\}}$ collects the ball- and pairing-weights of levels
$v-1,v$ and the pairings linking them to level $v-2$, $\qthreewt_{\phi(v)}$ collects
those of the single image level, and the second equality uses
$\kappa'_{\phi(v)}=\lambda'_{v-1}$. Two features make \eqref{eq:mergewt} strictly
stronger than the gap-decrement case and obstruct a direct appeal to the single-row
rule:
\begin{itemize}
\item the re-expansion of one image level into the two levels $v-1,v$ is a $1\to2$
operation, whereas the single-row Pieri/row identity \eqref{eq:Prow},
\cite[Lem.~5.4]{AMW25}, \cite[Lem.~5.6]{BDW25a} inserts \emph{one} new level over a
\emph{fixed} lower configuration and yields a single factor $e^*_m$, not a ratio of
the form $e^*_{\lambda'_{v-1}}e^*_{\lambda'_v}/e^*_{\kappa'_{\phi(v)}}$;
\item a column that read $b>a$ across the two merged values in $\lambda$ becomes a
$b=a$ column in $\kappa$, so by (W1) its ball weight changes from $x_k$ to
$x_k-t^{-(n-1)}$. Thus the lower level $v-1$ is itself \emph{modified} by the merge,
and one cannot hold it fixed while inserting level $v$, which is precisely the
hypothesis of the single-row rule.
\end{itemize}
We have not found a weight-preserving bijection or a Pieri-type identity that resolves
\eqref{eq:mergewt} accounting for these case changes, and we do not assert one. We
record \eqref{eq:mergewt} as the exact missing ingredient; numerically the ordinary
($x_k$-only) analogue is the Ayyer--Martin--Williams projection \eqref{eq:AMWproj},
which is a theorem \cite[Prop.~6.1]{AMW25}, making the interpolation version
\eqref{eq:mergewt} plausible, but a proof of \eqref{eq:mergewt} is beyond what we
establish here.

\emph{Conclusion.} By (A), (B), and (C): \eqref{eq:fiberwt} holds for every $\phi$
whose factorization in (A) uses only move \textup{(i)}; for general $\phi$ it is
conditional on the merge identity \eqref{eq:mergewt}. In particular the global decrement
$\phi(i)=\max(i-1,0)$ of Theorem~\ref{thm:G}, applied to a restricted $\lambda$ with
value set $U\subseteq\{0,2,3,\dots\}$, factors into moves \textup{(i)} only: since
$1\notin U$, every present value $v\ge2$ moves into the empty slot $v-1$ (a
gap-decrement) and $0$ is fixed, so the procedure of \textup{(A)} never invokes move
\textup{(ii)} (Lemma~\ref{lem:strict}). Hence Theorem~\ref{thm:G} is unconditional,
which is all that Theorem~\ref{thm:main} requires.
\end{proof}

\begin{qthreeremark}
The proof of Theorem~\ref{thm:main} invokes only Theorem~\ref{thm:G}, via
Corollary~\ref{cor:rec}. Theorem~\ref{thm:G} is exactly the decrement case and is
proved unconditionally by the complete weight-preserving bijection of
Lemma~\ref{lem:decrement}: there \emph{no} column case ever changes
(Lemma~\ref{lem:strict}), so the step does not depend on the merge resolution
\textup{(C)} of Proposition~\ref{prop:general}. Nothing in
Theorem~\ref{thm:main} uses the general (merge) case, which is why its conditional
status in Proposition~\ref{prop:general}\,\textup{(C)} is harmless here.
\end{qthreeremark}

\begin{qthreecorollary}\label{cor:rec}
Let $\rho$ be a composition with $\rho_i\ne1$ for all $i$, with $k$ nonzero parts,
and set $\eta=\rho^-$ (where $\rho^-_i=\max(\rho_i-1,0)$). Then at $q=1$,
\[
F^*_\rho(x;1,t)=F^*_\eta(x;1,t)\cdot e^*_k(x;t).
\]
\end{qthreecorollary}

\begin{proof}
Let $\lambda,\kappa$ reorder $\rho,\eta$. Take $\phi:i\mapsto\max(i-1,0)$, so
$\phi(\rho)=\eta$. \emph{Because $\lambda$ has no part of size $1$}, $\phi$ is a
bijection from $\{0,2,3,\dots\}$ onto $\{0,1,2,\dots\}$ (Lemma~\ref{lem:strict}),
so $\rho$ is the unique composition in $\qthreeSnl$ with $\phi(\rho)=\eta$; hence
$G^*_\eta=F^*_\rho$. By Lemma~\ref{lem:decrement} (equivalently
Theorem~\ref{thm:G}), $F^*_\rho=e^*_k\,F^*_\eta$, which is the claim; equivalently
$F^*_\rho/P^*_\lambda=F^*_\eta/P^*_\kappa$ together with
$P^*_\lambda/P^*_\kappa=e^*_k$ by \eqref{eq:factor} and the fact that $\kappa$ is
obtained from $\lambda$ by deleting the largest column (of size $k$).
\end{proof}

This is exactly where the restriction on $\lambda$ is essential: \emph{no part of
size $1$} makes $\phi$ strictly order-preserving on the occurring labels
(Lemma~\ref{lem:strict}), so the top-row deletion is a clean weight-preserving
bijection and the recursion closes; the \emph{unique part of size $0$} guarantees
$\eta$ has a unique $0$-part so the column-deletion is unambiguous. The hypotheses
are structural, not (as one might guess) a device to avoid a $0/0$ normalization.

\begin{qthreeremark}[two certified instances]
The dichotomy is visible already at $n=3$. For the restricted partition
$\lambda=(3,2,0)$ the decrement $\phi$ sends $3\mapsto2,\ 2\mapsto1,\ 0\mapsto0$,
which is strictly order-preserving on $\{0,2,3\}$: every column comparison
$(a,b)$ keeps its case ($=,>,<$) under $\phi$, so $\Psi$ of
Lemma~\ref{lem:decrement} is the weight-preserving bottom-row deletion
of \eqref{eq:rowshift} (the deleted row $1$ has $\lambda'_1=k=2$ balls), and
$F^*_\rho=e^*_2 F^*_{\rho^-}$ for each of the six
$\rho\in S_3(\lambda)$. For the \emph{non}-restricted $\lambda=(2,1,0)$ (a part of
size $1$) the decrement instead \emph{merges} the values $1$ and $0$
($\phi(1)=\phi(0)=0$): the column comparison $(1,0)$ changes from $>$ to $=$, so a
ball weight changes from $x_k$ to $x_k-t^{-(n-1)}$. This is precisely a merge move
\textup{(ii)} in the sense of Proposition~\ref{prop:general}\,\textup{(C)}, whose
fiber identity \eqref{eq:mergewt} we do \emph{not} prove; it is also why
Corollary~\ref{cor:rec} requires $\lambda$ to have no part of size $1$, so that the
recursion only ever uses gap-decrement moves and never a merge.
\end{qthreeremark}

\subsection*{From two-line queues to the dynamics}

\begin{qthreeproposition}\label{prop:step2}
For $\rho,\nu\in\qthreeSnl$ with $\rho_j=0$,
\[
P^{(2)}_j(\rho,\nu)=\frac{c^\rho_\nu}{\prod_{k<j}(x_k-t^{-n+2})\prod_{k>j}(x_k-t^{-n+1})},
\quad\text{equivalently}\quad
P_j\,P^{(2)}_j(\rho,\nu)=\frac{c^\rho_\nu}{e^*_{n-1}} .
\]
\end{qthreeproposition}

\begin{proof}
Write $D:=\prod_{k<j}(x_k-t^{-n+2})\prod_{k>j}(x_k-t^{-n+1})$, so that $P_jD=
c^\rho_\nu/e^*_{n-1}$ is the second assertion once the first is proved. Since
$\rho$ rearranges $\lambda$, which has distinct parts, $\mathcal G^\rho_\nu$ is empty
or a single element $\overline Q$ (Lemma~\ref{lem:ckm}), and $c^\rho_\nu=\qthreewt(\overline Q)$
in the latter case, $c^\rho_\nu=0$ in the former. We prove
\begin{equation}\label{eq:goal}
\qthreewt(\overline Q)=D\cdot P^{(2)}_j(\rho,\nu)\qquad(\overline Q\in\mathcal G^\rho_\nu),
\qquad\text{and}\qquad P^{(2)}_j(\rho,\nu)=0\ \text{when}\ \mathcal G^\rho_\nu=\varnothing,
\end{equation}
which is \eqref{eq:goal} in both cases.

\medskip
\noindent\emph{Step A: the dynamics is a tree of skip/settle decisions.}
In Step 2 a single active label $a$ is carried clockwise across the sites
$k=1,2,\dots,j-1$ (sites $k>j$ are never visited; site $j$ is the terminal site).
Initially $a=0$. On reaching a site $k<j$ holding a resident $b$ the chain makes
one of two moves:
\begin{itemize}
\item[(\emph{skip})] the resident $b$ stays at site $k$ and the \emph{same} label
$a$ continues to site $k+1$; the probability is $1-p_k$ if $b\ge a$ and $1-q_k$
if $b<a$;
\item[(\emph{settle})] the label $a$ stops at site $k$ (so the new content of
site $k$ is $a$), the resident $b$ becomes the new active label and continues to
site $k+1$; the probability is $p_k$ if $b\ge a$ and $q_k$ if $b<a$.
\end{itemize}
At site $j$ (whose original content is $\rho_j=0$) the current active label
settles, ending the step. Thus every realisation of Step 2 is a path in a binary
tree whose nodes are the visited sites $1,\dots,j-1$; the leaf at site $j$
determines the output configuration $\nu$, and
$P^{(2)}_j(\rho,\nu)$ is the sum over all root-to-leaf paths producing $\nu$ of
the product of the chosen branch probabilities.

\medskip
\noindent\emph{Step B: the weight emitted on each branch, divided by the matching
factor of $D$, is the branch probability.}
We assign to each branch a weight that is a single ball weight (and, for
settles, a factor $1-t$), and show that the emitted weight divided by the
corresponding factor of $D$ is exactly the branch probability listed in Step~A.
There are four branches, governed by the resident relation $b\ge a$ or $b<a$.
The matching factor of $D$ at a visited site $k<j$ is $x_k-t^{-n+2}$.

\smallskip
\noindent(1) \emph{skip, $b\ge a$.} The resident $b$ remains at site $k$, so in
the two-line queue the column-$k$ top ball (content of site $k$ \emph{before} the
move) and bottom ball (content \emph{after}) carry the same label $b$, i.e.\ the
case ``$b=a$'' of Definition~\ref{def:unsigned}; its ball weight is
$x_k-t^{-(n-1)}=x_k-t^{-n+1}$. Dividing by $x_k-t^{-n+2}$ gives, by the first
identity in \eqref{eq:skipprobs}, exactly $1-p_k$. We therefore emit the ball
weight $x_k-t^{-n+1}$ on this branch.

\smallskip
\noindent(2) \emph{skip, $b<a$.} Here the moving label $a$ passes over a strictly
weaker resident $b$. The resident again stays put (same-label column, ball weight
$x_k-t^{-n+1}$ as in (1)), but the moving label has jumped one weaker ball,
contributing one factor $t$ to the eventual pairing weight $(1-t)t^{\mathrm{skip}}$
of the pairing carrying $a$. Hence the emitted weight is $t\,(x_k-t^{-n+1})$, and
dividing by $x_k-t^{-n+2}$ gives, by the second identity in \eqref{eq:skipprobs},
exactly $1-q_k$.

\smallskip
\noindent(3) \emph{settle, $b\ge a$.} The label $a$ stops at site $k$ and the
resident $b\ge a$ is displaced. In the queue this opens a nontrivial pairing
that begins at column $k$ (the pairing that will carry $b$ onward); the new
column-$k$ ball relation is the displaced label below a smaller settled label,
which is the case ``$b<a$'' of Definition~\ref{def:unsigned}, of ball weight
$t^{-(n-1)}=t^{-n+1}$. The pairing just opened contributes its base factor
$1-t$ (its $t^{\mathrm{skip}}$ accumulates on later skip-branches of type (2)).
Thus the emitted weight is $t^{-n+1}(1-t)$, and dividing by $x_k-t^{-n+2}$ gives
exactly $p_k=t^{-n+1}(1-t)/(x_k-t^{-n+2})$.

\smallskip
\noindent(4) \emph{settle, $b<a$.} The label $a$ stops at site $k$ and the
strictly weaker resident $b<a$ is displaced; the new column-$k$ ball relation is
the displaced label below a larger settled label, the case ``$b>a$'' of
Definition~\ref{def:unsigned}, of ball weight $x_k$. As in (3) a nontrivial
pairing opens, contributing $1-t$. The emitted weight is $x_k(1-t)$, and dividing
by $x_k-t^{-n+2}$ gives exactly $q_k=(1-t)x_k/(x_k-t^{-n+2})$.

\smallskip
In all four cases the emitted weight at a visited site $k<j$, divided by the
factor $x_k-t^{-n+2}$ of $D$, equals the branch probability.

\medskip
\noindent\emph{Step C: the terminal site and the sites $k>j$.}
At the terminal site $j$ the active label settles; the original content there is
$\rho_j=0$, so column $j$ has a top vacancy and emits no ball weight (and is not a
factor of $D$). The sites $k>j$ are never visited by the sweep, so their content
is unchanged: $\nu_k=\rho_k$. In the queue, column $k>j$ then has equal top and
bottom labels, the case ``$b=a$'', of ball weight $x_k-t^{-(n-1)}=x_k-t^{-n+1}$,
which is precisely the matching factor of $D$ for $k>j$. Hence the column-$k$
contribution divided by its $D$-factor equals $1$, consistent with the statement
that columns $k>j$ never move.

\medskip
\noindent\emph{Step D: telescoping a settle-move into one pairing.}
A nontrivial pairing in $\overline Q$ corresponds to a maximal run in which a label
is carried forward: it is opened by a settle-branch (type (3) or (4)), then
crosses some number $s\ge0$ of weaker residents via skip-branches of type (2)
(each contributing one factor $t$), and finally settles again or reaches site
$j$. The factors accumulated on this run are $(1-t)$ from the opening settle and
$t$ from each of the $s$ type-(2) skips it traverses, giving total pairing weight
$(1-t)t^{s}=(1-t)t^{\mathrm{skip}(p)}$, in agreement with
Definition~\ref{def:unsigned}; here $\mathrm{skip}(p)=s$ is the number of weaker
balls the pairing jumps over. Each type-(1) skip and each unvisited site $k>j$
contributes its own same-label ball weight $x_k-t^{-n+1}$ as in (1) and Step~C.
Multiplying the contributions of all columns therefore yields exactly
$\qthreewt(\overline Q)$ on the one hand (product of the ball weights of all $n$
columns times $(1-t)t^{\mathrm{skip}}$ over all nontrivial pairings) and, on the
other hand, the product of all branch probabilities of the realising path times
$D$ (each visited factor $x_k-t^{-n+2}$ for $k<j$ supplied by Step~B and each
factor $x_k-t^{-n+1}$ for $k>j$ supplied by Step~C). That is, for a single
root-to-leaf path $\sigma$ realising $\nu$,
\[
\qthreewt_\sigma(\overline Q)=D\cdot(\text{product of branch probabilities along }\sigma).
\]

\medskip
\noindent\emph{Step E: uniqueness of the realising path and the empty case.}
A clockwise sweep moves each label only from its site toward site $j$ (a settle
fixes the active label at the current site and hands the resident forward) or
leaves it fixed (a skip). Consequently the output $\nu$ determines the path
$\sigma$ uniquely: reading $\nu$ from site $1$ to site $j$, the label $\nu_k$ at a
visited site $k<j$ was either the resident that stayed (skip at $k$) or the active
label that settled (settle at $k$), and which of the two occurred is forced by
whether $\nu_k=\rho_k$ or not, since $\lambda$ has distinct parts so the resident
and the active label are never equal except in a skip with $\nu_k=\rho_k=b$.
Thus there is \emph{exactly one} root-to-leaf path $\sigma$ producing a given
realisable $\nu$. Moreover the column ball-cases read off from $\sigma$ in
Steps~B--C are exactly the column ball-cases of the unique
$\overline Q\in\mathcal G^\rho_\nu$ (Lemma~\ref{lem:ckm}), and its nontrivial pairings
are exactly the settle-runs of $\sigma$ described in Step~D; hence
$\qthreewt_\sigma(\overline Q)=\qthreewt(\overline Q)=c^\rho_\nu$. Dividing the identity
$\qthreewt_\sigma(\overline Q)=D\cdot(\text{branch product along }\sigma)$ of Step~D
by $D$,
\[
P^{(2)}_j(\rho,\nu)=(\text{branch product along }\sigma)
=\frac{\qthreewt(\overline Q)}{D}=\frac{c^\rho_\nu}{D},
\]
which is \eqref{eq:goal}. Finally, if no two-line queue has top $\rho$ and bottom
$\nu$ (so $\mathcal G^\rho_\nu=\varnothing$ and $c^\rho_\nu=0$), then no sequence of
skip/settle moves transforms $\rho$ into $\nu$, because the realisable images of
$\rho$ under the sweep are precisely the bottom rows $\nu$ admitting such a queue
(each settle-run is a forward transport of a label, which is exactly a nontrivial
pairing, and the fixed labels are exactly the trivial same-label columns); hence
$P^{(2)}_j(\rho,\nu)=0=c^\rho_\nu/D$. This proves \eqref{eq:goal} in all cases,
and with it the proposition.

Finally, summing \eqref{eq:goal} over all $\nu$ (with $D$ fixed) and using that
the branch probabilities at each visited site sum to $1$
($p_k+(1-p_k)=q_k+(1-q_k)=1$) shows $\sum_\nu P^{(2)}_j(\rho,\nu)=1$, so
$P^{(2)}_j(\rho,\cdot)$ is genuinely a probability distribution; combined with
$0<t<1$ and $x_i>t^{-n+1}$ (whence $p_k,q_k\in[0,1]$ by \eqref{eq:skipprobs} and
positivity of the numerators), all branch probabilities lie in $[0,1]$.
\end{proof}

\begin{qthreeproposition}[Step 1 is a stochastic push onto a vacancy at site $j$]\label{prop:step1}
Fix a state $\mu\in\qthreeSnl$ and a site $j$, and run the Push-TASEP move of
Definition~\ref{def:chain}\,\textup{(Step 1)} starting by activating the label
$\mu_j$ at site $j$. Then:
\begin{enumerate}
\item[\textup{(a)}] every realisation of Step 1 terminates after finitely many
displacements;
\item[\textup{(b)}] every output configuration $\rho$ is again an element of
$\qthreeSnl$ (a rearrangement of $\lambda$) with $\rho_j=0$;
\item[\textup{(c)}] $P^{(1)}_j(\mu,\cdot)$ is a probability distribution on
$\{\rho\in\qthreeSnl:\rho_j=0\}$, i.e.\ $P^{(1)}_j(\mu,\rho)\ge0$ for all $\rho$
and $\sum_{\rho:\rho_j=0}P^{(1)}_j(\mu,\rho)=1$.
\end{enumerate}
\end{qthreeproposition}

\begin{proof}
A realisation of Step 1 is determined by the sequence of distinct sites it visits.
Write $s_0:=j$ and $b_0:=\mu_{s_0}=\mu_j$ for the activated label. If $b_0=0$
(the activated particle is the vacancy) there are no particles weaker than label
$0$, the cascade is empty, and the output is $\rho=\mu$ with $\rho_j=\mu_j=0$;
all three claims are trivial in this case. Assume henceforth $b_0>0$.

\smallskip
\noindent\emph{The displacement chain.} Suppose inductively that the active label
$b_{i-1}>0$ sits at site $s_{i-1}$. By definition of the move it jumps clockwise
to a \emph{weaker} particle, i.e.\ to a site whose current label is strictly
smaller than $b_{i-1}$; among the $m_i\ge1$ such weaker sites (in clockwise order
from $s_{i-1}$), it chooses the $k$-th with probability $t^{k-1}/[m_i]_t$, landing
on a new site $s_i$ that does not lie in $\{s_0,\dots,s_{i-1}\}$ (it is a fresh
weaker site). Set $b_i:=$ (label originally at $s_i$). Crucially the chosen target
is \emph{strictly weaker}, so
\begin{equation}\label{eq:strictdec}
b_0>b_1>b_2>\cdots\ge0 ,
\end{equation}
a strictly decreasing sequence of nonnegative integers. The move \emph{repeats}
exactly while the displaced label $b_i$ is positive, and \emph{ends} as soon as
$b_i=0$ (the active label reaches a vacancy).

\smallskip
\noindent\emph{(a) Termination, and that a weaker site always exists until the
vacancy is reached.} Since $\lambda$ is restricted it has a \emph{unique} part
equal to $0$, so $\mu$ has a unique vacancy, at some site $v$. As long as the
active label $b_{i-1}$ is positive, the vacancy is a particle (label $0$) strictly
weaker than $b_{i-1}$; it has not yet been visited (the chain ends the instant it
lands on a vacancy, and there is only one), so it is among the available weaker
sites and $m_i\ge1$. Hence the move is always defined while $b_{i-1}>0$. By the
strict decrease \eqref{eq:strictdec} the active label drops by at least $1$ at
each displacement, so after at most $b_0\le\lambda_1$ steps it reaches $0$ and the
cascade stops; in particular it stops precisely when it lands on the unique
vacancy $v$, say $s_r=v$, $b_r=0$.

\smallskip
\noindent\emph{(b) The output is in $\qthreeSnl$ with $\rho_j=0$.} The net effect
of the cascade is the cyclic-type relabelling along the visited chain
$s_0,s_1,\dots,s_r$: each visited label is carried one step forward along the
chain,
\[
\rho_{s_i}=b_{i-1}=\mu_{s_{i-1}}\quad(1\le i\le r),\qquad \rho_{s_0}=0,
\]
and $\rho_s=\mu_s$ for every site $s\notin\{s_0,\dots,s_r\}$. This is a bijection
of the multiset of labels on the ring: the labels $b_0,b_1,\dots,b_{r-1}$ occupy
the sites $s_1,\dots,s_r$, the consumed vacancy $b_r=0$ is removed from $s_r$ and
re-created at $s_0=j$, and all other sites are untouched. Hence the multiset of
labels of $\rho$ equals that of $\mu$, i.e.\ $\rho\in\qthreeSnl$, and by
construction $\rho_{s_0}=\rho_j=0$, as claimed. (Equivalently: the unique vacancy
has migrated from site $v$ to site $j$.)

\smallskip
\noindent\emph{(c) Stochasticity.} Nonnegativity is clear since each branch
probability $t^{k-1}/[m_i]_t$ is a quotient of nonnegative quantities for
$0<t<1$. For the row sum, observe that a realisation of Step 1 is a root-to-leaf
path in a finite decision tree: the root is the activation at $s_0=j$, the
children of a node with active label $b_{i-1}>0$ are its $m_i$ admissible weaker
targets $s_i$ (one per $k=1,\dots,m_i$), and a node is a leaf exactly when the
landed label is $0$. By part (a) every branch is finite. At each internal node the
outgoing branch probabilities sum to
\[
\sum_{k=1}^{m_i}\frac{t^{k-1}}{[m_i]_t}
=\frac{1+t+\cdots+t^{m_i-1}}{[m_i]_t}=\frac{[m_i]_t}{[m_i]_t}=1 .
\]
Therefore the total weight of all root-to-leaf paths is $1$: formally, summing the
path-weights from the leaves upward, each internal node contributes a factor
$\sum_k t^{k-1}/[m_i]_t=1$ times the (inductively unit) total weight of its
subtrees. Since distinct leaves give distinct visited chains and hence (by the
relabelling formula in (b)) distinct outputs $\rho$, grouping the leaves by their
output yields
\[
\sum_{\rho:\rho_j=0}P^{(1)}_j(\mu,\rho)
=\sum_{\text{leaves }\ell}\Pr(\ell)=1 .
\]
This proves (c). (The grouping is in fact one leaf per realisable $\rho$, but we
only need the total.)
\end{proof}

\begin{qthreeproposition}\label{prop:transition}
If $\lambda$ is restricted and $\mu,\nu\in\qthreeSnl$, then
$\displaystyle P(\mu,\nu)=\sum_{\rho\in\qthreeSnl}\frac{a^\mu_\rho\,c^\rho_\nu}{e^*_{n-1}}$.
\end{qthreeproposition}

\begin{proof}
Each $\rho\in\qthreeSnl$ has a unique vacancy, at the site $j_0(\rho)$ with
$\rho_{j_0(\rho)}=0$; by Proposition~\ref{prop:step1}\,(b) a Step-1 move that rings
the bell at site $j$ can produce $\rho$ only if $j=j_0(\rho)$, so for each $\rho$
exactly one value of $j$ contributes to the inner sum. The encoding of Step 1 by
classical two-line queues \cite[Lem.~5.4]{AMW25} gives, for $j=j_0(\rho)$,
\[
P^{(1)}_j(\mu,\rho)=a^\mu_\rho ,
\]
so that, combining with Proposition~\ref{prop:step2}
($P_j\,P^{(2)}_j(\rho,\nu)=c^\rho_\nu/e^*_{n-1}$),
\[
P(\mu,\nu)=\sum_j P_j\!\!\sum_{\rho:\rho_j=0}\!P^{(1)}_j(\mu,\rho)P^{(2)}_j(\rho,\nu)
=\sum_j\sum_{\rho:\rho_j=0}P^{(1)}_j(\mu,\rho)\,\frac{c^\rho_\nu}{e^*_{n-1}}
=\sum_{\rho\in\qthreeSnl}\frac{a^\mu_\rho c^\rho_\nu}{e^*_{n-1}}.
\]
\end{proof}

\begin{qthreecorollary}[$P$ is a stochastic matrix]\label{cor:stochastic}
For $\lambda$ restricted and $0<t<1$, $x_k>t^{-n+1}$ for all $k$, the matrix
$P=(P(\mu,\nu))_{\mu,\nu\in\qthreeSnl}$ is row-stochastic:
$P(\mu,\nu)\ge0$ for all $\mu,\nu$, and $\sum_{\nu\in\qthreeSnl}P(\mu,\nu)=1$ for
every $\mu$.
\end{qthreecorollary}

\begin{proof}
\emph{Nonnegativity.} Under the standing assumptions every factor occurring in
$P(\mu,\nu)$ is nonnegative. Indeed each bell weight $P_j$ is a ratio of products
of strictly positive numbers $x_k-t^{-n+2}>0$ and $x_k-t^{-n+1}>0$ (positive
because $x_k>t^{-n+1}>t^{-n+2}$ for $0<t<1$), hence $P_j\ge0$ and in fact
$P_j>0$; each Step-1 probability $P^{(1)}_j(\mu,\rho)\ge0$ by
Proposition~\ref{prop:step1}\,(c); and each Step-2 probability
$P^{(2)}_j(\rho,\nu)\ge0$ because, by \eqref{eq:skipprobs} and the positivity of
the numerators, the branch probabilities $p_k,q_k,1-p_k,1-q_k$ all lie in $[0,1]$
(this is the row-sum/positivity statement at the end of the proof of
Proposition~\ref{prop:step2}). As a finite sum of products of nonnegative numbers,
$P(\mu,\nu)\ge0$.

\smallskip
\emph{Row sums.} We sum the three-stage decomposition
\[
P(\mu,\nu)=\sum_{j=1}^n P_j\sum_{\rho\in\qthreeSnl:\rho_j=0}
P^{(1)}_j(\mu,\rho)\,P^{(2)}_j(\rho,\nu)
\]
over $\nu\in\qthreeSnl$ and evaluate stage by stage, innermost first. For fixed
$j$ and fixed $\rho$ with $\rho_j=0$, Step 2 is a probability distribution on
$\qthreeSnl$ by Proposition~\ref{prop:step2} (its closing paragraph):
\[
\sum_{\nu\in\qthreeSnl}P^{(2)}_j(\rho,\nu)=1 .
\]
Hence
\[
\sum_{\nu}P(\mu,\nu)
=\sum_{j=1}^n P_j\sum_{\rho:\rho_j=0}P^{(1)}_j(\mu,\rho)
 \underbrace{\sum_{\nu}P^{(2)}_j(\rho,\nu)}_{=1}
=\sum_{j=1}^n P_j\sum_{\rho:\rho_j=0}P^{(1)}_j(\mu,\rho) .
\]
Next, for each fixed $j$, Step 1 is a probability distribution supported on
$\{\rho:\rho_j=0\}$ by Proposition~\ref{prop:step1}\,(c):
\[
\sum_{\rho:\rho_j=0}P^{(1)}_j(\mu,\rho)=1 .
\]
Therefore
\[
\sum_{\nu}P(\mu,\nu)=\sum_{j=1}^n P_j .
\]
Finally the bell weights sum to $1$ by their very definition: with
$D_j:=\prod_{k<j}(x_k-t^{-n+2})\prod_{k>j}(x_k-t^{-n+1})$ we have
$P_j=D_j/e^*_{n-1}$ and, by the displayed formula for $e^*_{n-1}$ in
Definition~\ref{def:chain}, $e^*_{n-1}=\sum_{j=1}^n D_j$, so
\[
\sum_{j=1}^n P_j=\frac{\sum_{j=1}^n D_j}{e^*_{n-1}}=\frac{e^*_{n-1}}{e^*_{n-1}}=1 .
\]
Combining the three evaluations gives $\sum_{\nu\in\qthreeSnl}P(\mu,\nu)=1$, as
claimed. (Equivalently, in the algebraic form of
Proposition~\ref{prop:transition}, this is the identity
$\sum_\nu\sum_\rho a^\mu_\rho c^\rho_\nu=e^*_{n-1}$.)
\end{proof}

\subsection*{Proof of Theorem~\ref{thm:main}}

\begin{proof}
Fix $\nu\in\qthreeSnl$. By \cite[Thm.~1.15, Lem.~5.6]{BDW25a},
\[
F^*_\nu(x;1,t)=\sum_{\eta\in\qthreeN^n}F^{*\eta}_\nu(x;t)\,F^*_{\eta^-}(x;1,t),
\qquad
F^{*\eta}_\nu(x;t)=\sum_{\kappa\in\qthreeN^n}a^\eta_\kappa\,c^\kappa_\nu .
\]
By Corollary~\ref{cor:rec} (using that $\eta$ has a unique $0$-part),
$F^*_{\eta^-}(x;1,t)=F^*_\eta(x;1,t)/e^*_{n-1}(x;t)$. Substituting,
\[
F^*_\nu(x;1,t)=\sum_{\eta\in\qthreeN^n}F^*_\eta(x;1,t)\sum_{\kappa\in\qthreeN^n}
\frac{a^\eta_\kappa\,c^\kappa_\nu}{e^*_{n-1}(x;t)}
=\sum_{\eta\in\qthreeN^n}F^*_\eta(x;1,t)\,P(\eta,\nu),
\]
the last equality by Proposition~\ref{prop:transition}. Thus the vector
$u:=\bigl(F^*_\mu(x;1,t)\bigr)_{\mu\in\qthreeSnl}$ is left-fixed by the transition
matrix $P$, i.e.\ $uP=u$.

It remains to deduce that $u$, suitably normalized, is \emph{the} (unique)
stationary distribution. This requires three facts about $P$, established in
Lemmas~\ref{lem:positivity}--\ref{lem:aperiodic} below: (i) $u$ is entrywise
strictly positive (Lemma~\ref{lem:positivity}), so that $\pi^*_\lambda:=u/\!\sum_\mu u_\mu$
is a genuine probability vector; (ii) $P$ is irreducible on $\qthreeSnl$
(Lemma~\ref{lem:irreducible}); and (iii) $P$ is aperiodic
(Lemma~\ref{lem:aperiodic}). Granting these, the standard Perron--Frobenius theory
of finite Markov chains applies (the matrix $P$ being row-stochastic by
Corollary~\ref{cor:stochastic}): an irreducible finite stochastic matrix has a
\emph{unique} stationary distribution, which is moreover strictly positive, and
every nonnegative left-fixed vector is a scalar multiple of it
\cite[Thm.~1.5.6, Thm.~1.7.7]{Norris}. Since $u$ is a strictly positive
left-fixed vector, it is therefore a positive multiple of the unique stationary
distribution; normalizing and using the symmetrization identity
$\sum_{\mu\in\qthreeSnl}F^*_\mu(x;1,t)=P^*_\lambda(x;1,t)$ (the $q=1$ case of the
fact that the interpolation ASEP polynomials of content $\lambda$ sum to the
interpolation Macdonald polynomial) to evaluate the normalizing constant gives
\[
\pi^*_\lambda(\mu)=\frac{u_\mu}{\sum_{\nu}u_\nu}
=\frac{F^*_\mu(x;1,t)}{P^*_\lambda(x;1,t)},\qquad \mu\in\qthreeSnl .
\]
By aperiodicity, $P^m(\eta,\cdot)\to\pi^*_\lambda$ as $m\to\infty$ from every
initial state $\eta$, so $\pi^*_\lambda$ is indeed the limiting/stationary law of
the chain.
\end{proof}

\subsection*{Positivity, irreducibility, and aperiodicity}

We work throughout under the standing assumptions $0<t<1$ and $x_k>t^{-n+1}$ for
all $k$, which are exactly those used in Definition~\ref{def:chain} to make all
branch probabilities lie in $[0,1]$. Recall that the only part of size $0$ is
unique, so every $\mu\in\qthreeSnl$ has a \emph{unique vacancy} site, the site
$j_0(\mu)$ with $\mu_{j_0}=0$.

\begin{qthreelemma}[strict positivity of the left eigenvector]\label{lem:positivity}
For every $\mu\in\qthreeSnl$ one has $F^*_\mu(x;1,t)>0$. Consequently
$P^*_\lambda(x;1,t)=\sum_{\mu}F^*_\mu(x;1,t)>0$ and
$\pi^*_\lambda(\mu)=F^*_\mu/P^*_\lambda\in(0,1)$ defines a probability
distribution on $\qthreeSnl$.
\end{qthreelemma}

\begin{proof}
By the combinatorial formula \eqref{eq:Fformula},
$F^*_\mu(x;1,t)=\sum_{Q:\,\mathrm{bot}(Q)=\mu}\qthreewt(Q)$, a sum over the signed
multiline queues of content $\lambda$ with bottom row $\mu$. This index set is
nonempty: the ``identity'' sMLQ in which the column of each label-value $v$
occupies rows $1,\dots,v$ over the site dictated by $\mu$, with all pairings
trivial, has bottom row $\mu$. Every summand $\qthreewt(Q)$ is a product of ball
weights and pairing weights. The ball weights are, by
Definition~\ref{def:unsigned}, one of
\[
x_k-t^{-(n-1)},\qquad x_k,\qquad t^{-(n-1)} ,
\]
each of which is strictly positive: $t^{-(n-1)}=t^{-n+1}>0$, and
$x_k-t^{-(n-1)}=x_k-t^{-n+1}>0$ by the standing assumption $x_k>t^{-n+1}$. The
pairing weights are $(1-t)\,t^{\mathrm{skip}}>0$ since $0<t<1$. Hence every
$\qthreewt(Q)>0$, and $F^*_\mu$ is a nonempty sum of positive terms, so
$F^*_\mu>0$. The remaining assertions are immediate.
\end{proof}

\begin{qthreelemma}[irreducibility]\label{lem:irreducible}
The interpolation $t$-Push TASEP is irreducible on $\qthreeSnl$: for any
$\mu,\nu\in\qthreeSnl$ there is $m\ge1$ with $P^m(\mu,\nu)>0$.
\end{qthreelemma}

\begin{proof}
It suffices to exhibit a family of \emph{single-step} transitions of strictly
positive probability whose induced moves act transitively on $\qthreeSnl$. We use
the following move, the \emph{vacancy swap}.

\smallskip
\noindent\emph{Vacancy swap $T_j$.} Fix $\mu$ with vacancy at site $p:=j_0(\mu)$,
and fix any site $j$ with $\mu_j=:a>0$. Ring the bell at site $j$ (probability
$P_j>0$, since $P_j$ is a ratio of strictly positive products by the standing
assumption). In Step~1 the active label $a>0$ travels clockwise; the vacancy at
$p$ is one of the weaker particles (label $0<a$), so among the admissible
push-targets the active may move \emph{directly to the vacancy site $p$}. This
choice has probability $t^{k-1}/[m]_t>0$, where the vacancy is the $k$-th weaker
particle clockwise from $j$ and $m\ge1$ is the total number of weaker particles
($[m]_t=1+t+\cdots+t^{m-1}>0$). Because the target $p$ holds a vacancy, the
cascade ends immediately: no intermediate site is altered, $a$ now sits at $p$,
and site $j$ becomes the new vacancy. Thus Step~1 outputs
$\rho=(j\;p)\!\cdot\!\mu$, the configuration obtained from $\mu$ by exchanging the
contents of sites $j$ and $p$ (so $\rho_p=a$, $\rho_j=0$, $\rho_k=\mu_k$
otherwise), and $\rho_j=0$ as required. In Step~2 the freshly created vacancy at
site $j$ sweeps from site $1$; choosing the \emph{skip} branch at every visited
site $1,\dots,j-1$ (each skip has probability $1-p_k>0$ or $1-q_k>0$ by
\eqref{eq:skipprobs} and the standing assumption) returns the vacancy to site $j$
and leaves $\rho$ unchanged. Hence
\[
P\bigl(\mu,\;(j\;p)\!\cdot\!\mu\bigr)\;\ge\;
P_j\cdot\frac{t^{\,k-1}}{[m]_t}\cdot\!\!\prod_{1\le k<j}\!(\text{skip prob at }k)\;>\;0 .
\]

\smallskip
\noindent\emph{Transitivity.} The moves $T_j$ realize, from any state $\mu$, the
transposition of the vacancy with the contents of \emph{any} site. We claim these
generate all of $\qthreeSnl$. Indeed, identify a configuration with the placement
of the $n-1$ distinct positive labels and one blank on the $n$ sites; a vacancy
swap exchanges the blank with a chosen labelled token. Given a target $\nu$,
sort $\mu$ into $\nu$ greedily: repeatedly move the blank (by a vacancy swap that
brings to the blank's site the label that $\nu$ requires there, then advancing)
one obtains $\nu$ in finitely many vacancy swaps. Concretely, transpositions of a
fixed moving point (the blank) with arbitrary points generate the full symmetric
group $S_n$ acting on the $n$ sites, and this action is transitive on the
arrangements with prescribed multiset of labels, i.e.\ on $\qthreeSnl$. Hence for
any $\mu,\nu$ there is a finite sequence of vacancy swaps carrying $\mu$ to $\nu$,
each of positive probability; multiplying, $P^m(\mu,\nu)>0$ for the corresponding
$m$. This proves irreducibility.
\end{proof}

\begin{qthreeremark}
Lemma~\ref{lem:irreducible} uses only the existence of \emph{some} positive-probability
single step for each vacancy swap; we do not need the exact transition probabilities,
only their positivity, which is guaranteed by the standing assumptions
$0<t<1,\ x_k>t^{-n+1}$. We verified the generation claim by exhaustive search for all
restricted $\lambda$ with $n\le5$, in agreement with the classical fact that
``swap-the-blank'' moves connect all token arrangements on $n$ sites.
\end{qthreeremark}

\begin{qthreelemma}[aperiodicity]\label{lem:aperiodic}
For every $\mu\in\qthreeSnl$ one has $P(\mu,\mu)>0$; hence the chain is aperiodic.
\end{qthreelemma}

\begin{proof}
Let $p:=j_0(\mu)$ be the vacancy site and ring the bell at $j:=p$ (probability
$P_p>0$). In Step~1 the activated ``particle'' at site $j=p$ carries label
$a=0$; there are $m=0$ particles weaker than a vacancy, so the cascade is empty
and Step~1 leaves the configuration unchanged, outputting $\rho=\mu$ with the
vacancy again at site $p$. In Step~2 the vacancy sweeps from site $1$; choosing
the \emph{skip} branch at each visited site $1,\dots,p-1$ (probability $1-p_k>0$
or $1-q_k>0$ at site $k$, by \eqref{eq:skipprobs} and the standing assumption)
returns the vacancy to site $p$ without moving any label, so the output is
$\nu=\mu$. Therefore
\[
P(\mu,\mu)\;\ge\;P_p\cdot\!\!\prod_{1\le k<p}\!(\text{skip prob at }k)\;>\;0 .
\]
A finite Markov chain with $P(\mu,\mu)>0$ for some state in each communicating
class is aperiodic; combined with irreducibility (Lemma~\ref{lem:irreducible})
the whole chain is aperiodic.
\end{proof}

\section*{Remark on the failure mode to avoid}

Replacing this chain by the ordinary multispecies ASEP (matrix-product ansatz /
Zamolodchikov--Faddeev algebra) does \emph{not} solve the stated problem: that
construction yields the ordinary Macdonald polynomials as stationary weights
\cite{CMW22,AMW25}, not the \emph{interpolation} polynomials $F^*_\mu,P^*_\lambda$.
The interpolation correction is precisely what the bell $P_j$ and the
push-with-sweep dynamics supply, and the proof above runs through signed
two-line queues rather than a trace over an auxiliary algebra.

\begin{thebibliography}{9}
\bibitem{AS19} P.~Alexandersson, M.~Sawhney.
  \emph{Properties of non-symmetric Macdonald polynomials at $q=1$ and $q=0$}.
  Ann.\ Comb.\ \textbf{23} (2019), 219--239.
\bibitem{AMW25} A.~Ayyer, J.~Martin, L.~Williams.
  \emph{The inhomogeneous $t$-PushTASEP and Macdonald polynomials at $q=1$}.
  Ann.\ Inst.\ Henri Poincar\'e D, 2025.
\bibitem{BDW25a} H.~Ben Dali, L.~Williams.
  \emph{A combinatorial formula for interpolation Macdonald polynomials}.
  Preprint, arXiv:2510.02587, 2025.
\bibitem{BDW26} H.~Ben Dali, L.~Williams.
  \emph{A probabilistic interpretation for interpolation Macdonald polynomials}.
  Preprint, arXiv:2602.13492, 2026.
\bibitem{Norris} J.~R.~Norris.
  \emph{Markov Chains}. Cambridge Series in Statistical and Probabilistic
  Mathematics. Cambridge University Press, 1997.
\bibitem{CMW22} S.~Corteel, O.~Mandelshtam, L.~Williams.
  \emph{From multiline queues to Macdonald polynomials via the exclusion process}.
  Amer.\ J.\ Math.\ \textbf{144} (2022), 395--436.
\bibitem{Dolega} M.~Dolega.
  \emph{Strong factorization property of Macdonald polynomials and higher-order
  Macdonald's positivity conjecture}.
  J.\ Algebraic Combin.\ \textbf{46} (2017), 135--163.
\end{thebibliography}

\endgroup
